\subsection{希尔伯特空间}

定义 5. --- 希尔伯特空间（或 Hilbert 空间）是 Hausdorff 且完备的准希尔伯特空间。我们说向量空间 E（K 上的）上的范数是希尔伯特的，如果它是准希尔伯特的，且赋范空间 E 是完备的。

若 E 是希尔伯特空间且 M 是 E 的闭向量子空间，则 M 上诱导的准希尔伯特空间结构实际上是希尔伯特空间结构。这时我们说赋有诱导结构的 M 是 E 的希尔伯特子空间。

例. --- 1) V, p.~4 的例 1、3、4 中定义的准希尔伯特空间都是希尔伯特空间。另一方面，例 2 中定义的准希尔伯特空间既非 Hausdorff，也非完备。希尔伯特空间的复化是希尔伯特空间。

* 2) 设 X 为 Hausdorff 拓扑空间，\(\mu\) 为 X 上的正测度。设 \(L^2(X, \mu)\) 为 X 上取值于 C 的全体平方 \(\mu\) -可积函数按 \(\mu\) 的等价类组成的空间。这是一个复希尔伯特空间，其内积由下式给出

\[
\langle f | g \rangle = \int_ {\mathrm{X}} \overline {{{f (x)}}} g (x) d \mu (x). *
\]

* 3) 设 \(n \geqslant 1\) 为整数，\(U\) 为 \(\mathbf{R}^n\) 中的开集。设 \(\mu\) 为 \(U\) 上由 \(\mathbf{R}^n\) 上的 Lebesgue 测度诱导的测度，令 \(\mathcal{H}^0 = L^2(U, \mu)\) 。用 \(\mathcal{H}^1\) 表示具有下列性质的所有函数 \(f \in \mathcal{H}^0\) 组成的空间；对 \(1 \leqslant i \leqslant n\) ，存在函数 \(g_i \in \mathcal{H}^0\) 使

\[
\int_ {\mathrm{U}} g _ {i} (x) h (x) d \mu (x) = - \int_ {\mathrm{U}} f (x) \mathrm{D} _ {i} h (x) d \mu (x)
\]

对每个函数 \(h\) （\(\mathbf{C}^1\) 类的，在 \(\mathbf{U}\) 中有紧支撑的）成立。函数 \(g_i\) 在相差一个关于 \(\mu\) 的等价类的意义下唯一确定，记作 \(\mathbf{D}_i f\) 或 \(\partial f / \partial x_i\) （第 i 个偏导数）。对整数 \(s \geqslant 1\) 用归纳法，把 \(\mathcal{H}^s\) 定义为所有函数 \(f \in \mathcal{H}^1\) （使 \(D_i f \in \mathcal{H}^{s-1}\) 对 \(1 \leqslant i \leqslant n\) 成立的）组成的集。我们用下列公式在 \(\mathcal{H}^s\) 上定义内积

\[
\langle f | g \rangle = \sum_ {k = 0} ^ {s} \sum_ {1 \leqslant i _ {1} \leqslant \dots \leqslant i _ {k} \leqslant n} \int \overline {{{{\mathrm{D} _ {i _ {1}} \dots \mathrm{D} _ {i _ {k}} f}}}}. \mathrm{D} _ {i _ {1}} \dots \mathrm{D} _ {i _ {k}} g d \mu .
\]

于是 \(\mathcal{H}^s\) 是一个复希尔伯特空间，称为指标 \(s\) 的 Sobolev 空间。

* 4) 设 X 为 \(\mathbf{C}^r\) 类（\(r \geqslant 1\) ）纯的、维数有限为 \(n\) 的微分流形。

在向量纤维空间 \(\Lambda^{n}\mathrm{T}(\mathrm{X})\) 中，设 L 为零截面的余集。对每个实数 \(\lambda \neq 0\) ，映射 \(u \mapsto \lambda u\) （从 \(\Lambda^{n}\mathrm{T}(\mathrm{X})\) 到自身的）使 L 稳定。

设 \(\alpha\) 为复数。一个复值函数 \(\omega\) （定义在 \(\mathbf{L}\) 上），若 \(\omega (\lambda u) = |\lambda |^{\alpha}\omega (u)\) 对 \(u\in \mathbf{L}\) 与任意非零实数 \(\lambda\) 成立，就称为 \(\alpha\) 阶密度（\(\mathbf{X}\) 上的）。我们说 1 阶密度 \(\omega\) 是局部可积的，如果存在一个开覆盖 \((\mathrm{U}_i)_{i\in \mathbf{I}}\) （\(\mathbf{X}\) 的），且对每个 \(i\in \mathbf{I}\) ，存在一个坐标系 \(\xi_i = (\xi_i^1,\dots,\xi_i^n)\) （\(\mathrm{U}_i\) 上的）与一个复值函数 \(f_{i}\) （\(\xi_i(\mathrm{U}_i)\) 上的），满足下列条件：a) 函数 \(f_{i}\) 在开集 \(\xi_i(\mathrm{U}_i)\) （\(\mathbf{R}^n\) 中的）上关于 Lebesgue 测度 \(\mu\) 是局部可积的；

\begin{enumerate}
\def\labelenumi{\alph{enumi})}
\setcounter{enumi}{1}
\tightlist
\item
  设 \(x \in \mathbf{U}_i\) ；若 \((\partial_{1,i,x}, \ldots, \partial_{n,i,x})\) 是 \(\mathrm{T}_x\mathrm{X}\) 的、与坐标系 \((\xi_i^1, \ldots, \xi_i^n)\) （\(\mathbf{U}_i\) 中的）相伴的基，则有
\end{enumerate}

\[
\omega \left(\partial_ {1, i, x} \wedge \dots \wedge \partial_ {n, i, x}\right) = f _ {i} \left(\xi_ {i} ^ {1} (x), \dots , \xi_ {i} ^ {n} (x)\right).
\]

于是，存在唯一的测度 \(\tilde{\omega}\) （\(\mathbf{X}\) 上的），使对每个 \(i\in \mathrm{I}\) ，在 \(\xi_{i}\) 下 \(\tilde{\omega}\) 在 \(\mathrm{U}_i\) 上的限制的像等于测度 \(f_{i}.\mu\) （参见~VAR, R, 10.4.3）。

设 \(\mathcal{V}\) （相应地，\(\mathcal{N}\) ）为由可测密度 \(\omega\) （\(1/2\) 阶）组成的、使与 1 阶密度 \(|\omega|^2\) 相伴的测度有界（相应地，为零）的向量空间。设 \(\omega_1\) 与 \(\omega_2\) 属于 \(\mathcal{V}\) ；则 \(\omega = \overline{\omega}_1\omega_2\) 是 1 阶密度，且测度 \(\tilde{\omega}\) （与 \(\omega\) 相伴的）有界；数 \(\int_{\mathbf{X}}\tilde{\omega}\) 只依赖于类 \(\dot{\omega}_1\) 与 \(\dot{\omega}_2\) （即 \(\omega_1\) 与 \(\omega_2\) 模 \(\mathcal{N}\) 的类），记作 \(\langle\omega_1|\omega_2\rangle\) 或 \(\langle\dot{\omega}_1|\dot{\omega}_2\rangle\) 。于是映射 \((\dot{\omega}_1,\dot{\omega}_2)\mapsto\langle\dot{\omega}_1|\dot{\omega}_2\rangle\) 给向量空间 \(\Omega_{1/2}(\mathbf{X})=\mathcal{V}/\mathcal{N}_{*}\) 赋以一个复希尔伯特空间结构。

*5) 设 D 为 C 中以 0 为心、1 为半径的开圆盘。Hardy 空间 \(\mathrm{H}^2 (\mathbf{D})\) 由满足下列条件的全纯函数 \(f\colon \mathbf{D}\to \mathbf{C}\) 组成：

\[
\sup _ {0 <   \mathbf {R} <   1} \int_ {0} ^ {1} | f (\mathbf {R}. \mathbf {e} (\theta)) | ^ {2} d \theta <   + \infty .
\]

若 \(f_{1}\) 与 \(f_{2}\) 属于 \(\mathrm{H}^2 (\mathbf{D})\) ，则极限

\[
\langle f _ {1} | f _ {2} \rangle = \lim _ {\mathbf {R} \rightarrow 1} \int_ {0} ^ {1} \overline {{{f _ {1} (\mathbf {R} . \mathbf {e} (\theta))}}}. f _ {2} (\mathbf {R}. \mathbf {e} (\theta)) d \theta
\]

存在；映射 \((f_1, f_2) \mapsto \langle f_1 | f_2 \rangle\) 给向量空间 \(\mathbf{H}^2(\mathbf{D})\) 赋以一个复希尔伯特空间结构。

为使函数 \(f: \mathbf{D} \to \mathbf{C}\) 属于 \(\mathbf{H}^2(\mathbf{D})\) ，必要且充分的是存在复数的序列 \((a_n)_{n \in \mathbb{N}}\) ，使 \(\sum_{n=0}^{\infty} |a_n|^2 < +\infty\) 且

\[
f (z) = \sum_ {n = 0} ^ {\infty} a _ {n} z ^ {n}
\]

对所有 \(z \in \mathbf{D}\) 成立。这时有 \(\| f\|^2 = \sum_{n=0}^{\infty} |a_n|^2\) ，由此给出从 \(\mathrm{H}^2(\mathrm{D})\) 到希尔伯特空间 \(\ell^2\) 上的一个同构（V, p.~4）。 *

每个 Hausdorff 准希尔伯特空间都同构于一个在相差一个同构的意义下确定的希尔伯特空间的处处稠密子空间。确切地说：

命题 4. --- 设 E 为 Hausdorff 准希尔伯特空间，\(\hat{E}\) 为 E 的赋范空间完备化（GT, IX, § 3, No.~3）。内积 \((x, y) \mapsto \langle x|y \rangle\) 由连续性延拓为 \(\hat{\mathbf{E}}\) 上的正且非退化的埃尔米特型，并在 \(\hat{\mathbf{E}}\) 上定义一个希尔伯特空间结构。

\((x, y) \mapsto \langle x|y \rangle\) 到 \(\hat{E} \times \hat{E}\) 上的延拓的存在性由这个半双线性型在 \(E \times E\) 上的连续性得出（GT, III, §6, No.~5, 定理 1）。此外，这个延拓（仍记作 \((x, y) \mapsto \langle x|y \rangle\) ）是埃尔米特型且满足关系 \(\langle x|x\rangle = \|x\|^2\) ，这由恒等式延拓原理得出（ \(\|x\|\) 是 \(\hat{E}\) 上把 E 上的范数用连续性延拓所得的范数）；这就证明了关系 \(\langle x|x\rangle = 0\) 蕴含在 \(\hat{E}\) 中 x = 0，从而型 \((x, y) \mapsto \langle x|y \rangle\) 是正的且非退化的，于是在 \(\hat{E}\) 上定义一个希尔伯特空间结构。

证毕。

这个希尔伯特空间称为 Hausdorff 准希尔伯特空间 E 的完备化。

* 例 6. --- 设 U 为 \(\mathbf{R}^n\) 的开子集（ \(n \geqslant 1\) ）。设 \(\mathcal{C}_0^1(\mathrm{U})\) 为 U 中具有紧支撑的 \(\mathbf{C}^1\) 类函数全体组成的向量空间。我们在 \(\mathcal{C}_0^1(\mathrm{U})\) 上定义一个 Hausdorff 准希尔伯特空间结构，其内积由下式给出

\[
\langle f | g \rangle = \sum_ {i = 1} ^ {n} \int_ {\mathrm{U}} \overline {{{{\mathbf {D} _ {i} f (x)}}}}. \mathbf {D} _ {i} g (x) d x.
\]

这个准希尔伯特空间不完备。它的完备化称为与 U 相伴的 Dirichlet 空间。 *

推论. --- 设 V 为 K 上的向量空间，f 为 V 上的正埃尔米特型。a) 存在一个 Hilbert 空间 E 和一个线性映射 \(u:V \rightarrow E\) ，使 \(f(x,y) = \langle u(x)|u(y) \rangle\) 对 V 中的 x, y 成立，且使 \(u(V)\) 在 E 中稠密。

\begin{enumerate}
\def\labelenumi{\alph{enumi})}
\setcounter{enumi}{1}
\tightlist
\item
  若两个偶 \((\mathbf{E}_{i}, u_{i})\) 满足与 a) 类似的条件，则存在唯一的同构 \(\phi\) （从 Hilbert 空间 \(E_{1}\) 到 Hilbert 空间 \(E_{2}\) 上的）使 \(u_{2} = \phi \circ u_{1}\) 。
\end{enumerate}

设 N 为全体 \(x \in V\) （使 \(f(x, x) = 0\) 的）组成的集。我们在空间 V/N 上用 \(\langle \dot{x} | \dot{y} \rangle = f(x, y)\) （对 \(x \in \dot{x}\) 与 \(y \in \dot{y}\) ）定义一个正且非退化的埃尔米特型。设 E 为 V/N 的希尔伯特空间完备化，u 为从 V 到 E 中的映射 \(x \mapsto x + N\) 。于是 a) 的条件满足。

在 b) 的假设下，N 等于 \(u_{1}\) 的核，也等于 \(u_{2}\) 的核。因此存在双射线性映射 \(\phi_{0}\) （从 \(u_{1}(V)\) 到 \(u_{2}(V)\) 上的），使 \(u_{2}(x) = \phi_{0}(u_{1}(x))\) 对所有 \(x \in V\) 成立。立即可以验证 \(\phi_{0}\) 是准希尔伯特空间之间的同构，从而是一个等距。由于对 i = 1, 2，\(u_{i}(V)\) 在 \(E_{i}\) 中稠密，\(\phi_{0}\) 唯一地延拓为一个等距 \(\phi\) （从 \(E_{1}\) 到 \(E_{2}\) 上的），由此得 b)。

我们说希尔伯特空间 \(\mathbf{E}\) 是 \(\mathbf{V}\) 的分离完备化（关于型 \(f\) ）。

例 7. --- 设 G 为一个群（单位元为 1），\(\pi\) 为从 G 到一个复希尔伯特空间 E 的自同构群中的同态；我们说 \(\pi\) 是 G 在 E 中的一个酉表示。设 \(a \in E\) ；我们令

\[
\phi (x) = \langle a | \pi (x). a \rangle
\]

对所有 \(x \in \mathbf{G}\) 成立。则 \(\phi : \mathbf{G} \to \mathbf{C}\) 是正定的，换言之满足关系：

(PD) 对 C 中每个 \(\lambda_{1},\ldots,\lambda_{n}\) 与 G 中每个 \(x_{1},\ldots,x_{n}\) ，有

\[
\sum_ {i, j = 1} ^ {n} \overline {{\lambda}} _ {i} \lambda_ {j} \phi (x _ {i} ^ {- 1} x _ {j}) \geqslant 0.\tag{11}
\]

事实上，(11) 的左端正是 \(\| \sum_{i=1}^{n} \lambda_i \pi(x_i). a \|^2\) 。

反之，设 \(\phi\) 为 \(G\) 上的正定函数。设 \(\mathbf{C}^{(G)}\) 为 \(G\) 上具有有限支撑的全体函数组成的向量空间。我们定义一个埃尔米特型 \(\Phi\) 于 \(G\) 上：

\[
\Phi (u, v) = \sum_ {x, y \in G} \overline {{u (x)}} v (y) \phi (x ^ {- 1} y)\tag{12}
\]

而关系式 (PD) 表明 \(\Phi\) 是正的。由命题 4 的推论，存在一个希尔伯特空间 \(E\) 和一个线性映射 \(\rho : \mathbf{C}^{(G)} \to E\) ，其像稠密，使得

\[
\Phi (u, v) = \langle \rho (u) | \rho (v) \rangle \quad \text { for } \quad u, v \quad \text { in } \quad \mathbf {C} ^ {(G)}.\tag{13}
\]

对每个 \(x \in \mathbf{G}\) ，设 \(\gamma_x\) 为 \(x\) 在 \(\mathbf{C}^{(\mathrm{G})}\) 中的左平移，由 \(\gamma_x u(y) = u(x^{-1}y)\) （对 \(u \in \mathbf{C}^{(\mathrm{G})}\) 与 \(y \in \mathbf{G}\) ）定义。我们有 \(\Phi(\gamma_x u, \gamma_x v) = \Phi(u, v)\) 。现在把命题 4 的推论的断言 \(b\) 应用于 \(\rho\) 与 \(\rho \circ \gamma_x\) ：存在唯一的自同构 \(\pi(x)\) （希尔伯特空间 \(E\) 上的）使 \(\rho \circ \gamma_x = \pi(x) \circ \rho\) 。立即可以看出 \(\pi\) 是从 \(G\) 到 \(E\) 的自同构群中的同态。

设 \(\delta\) 为 \(\mathbf{C}^{(\mathrm{G})}\) 中由 \(\delta(1) = 1\) 与 \(\delta(x) = 0\) （对 \(x \neq 1\) ，\(\mathbf{G}\) 中的）定义的元素。我们有 \(u = \sum_{x \in \mathbf{G}} u(x) \cdot \gamma_x \delta\) 对所有 \(u \in \mathbf{C}^{(\mathrm{G})}\) 成立，于是 \(\rho(u) = \sum_{x \in \mathbf{G}} u(x) \pi(x) \cdot a\) （令 \(a = \rho(\delta)\) ）。公式 (12) 与 (13) 蕴含 \(\phi(x) = \langle a | \pi(x) \cdot a \rangle\) 对所有 \(x \in \mathbf{G}\) 成立。我们注意，向量 \(\pi(x) \cdot a\) （对所有 \(x \in \mathbf{G}\) ）组成的集在 \(\mathbf{E}\) 中是完全的。

